\section{Benchmark design}
Figure~\ref{fig:ladder} summarizes a conceptual ladder rather than a claim that every historical stage is final evidence. Stage I demonstrates that ordinary distributions can be shortcut-heavy. Stage II matches static information. Stage III additionally matches unordered history statistics and changes signed temporal direction. Stage IV constructs nonlinear latent-dynamics pairs with common final two frames. Stage V uses shared but action-diverse candidate sets and independent physical-parameter, temporal-delay, and held-mechanism OOD cells.

Each label is obtained by a GPU-PhysX proxy-to-target execution under the same threshold and horizon, not by the analytic baseline formula. The full v2 configuration fixes six history frames, a tolerance of $10^{-6}$, the observable pair-swap, candidate counts C5/C10/C20, grouped splits, and an execution budget of 92,000. This separation matters: analytic baselines estimate from fair observations, while execution determines success.

\begin{figure}[t]\centering\includegraphics[width=\linewidth]{shortcut_removal_layers.png}\caption{Shortcut-removal ladder. The figure is conceptual; quantitative values remain in frozen 3S tables. Source: \texttt{milestone3s\_paper\_package/figures/figure\_provenance.json}.}\label{fig:ladder}\end{figure}
\begin{figure}[t]\centering\includegraphics[width=\linewidth]{counterfactual_pair_construction.png}\caption{Final-pair construction: early frames are swapped while common final synchronization frames and the candidate set are retained. This is the valid operator used for the causal audit, unlike generic full reversal. Source: frozen v2 probe configuration and forensic audit.}\label{fig:pair}\end{figure}
